% 【本文件写什么】impl 实现层：dummy
%   - 定位：占位/链接桩实现，保证默认 feature 可编译。
%   - crate 路径：os/components/wateros-runtime/runtime-console/console-impl/impl-dummy/src/
%   - 如何实现对应 api-v0 trait：关键类型、状态机、数据结构
%   - 算法或硬件交互流程（可用伪代码/流程图）
%   - 本 impl 启用的 Cargo [features] 与 arch feature（impl-riscv64 等）
%   - 与其它 impl 变体的差异、切换 feature 时的影响
%   - 性能/内存假设与已知限制
% 【事实来源】
%   - 上述 crate 源码与 Cargo.toml
%   - os/feature-tree.txt 中指向本 impl 的 feature 链

\subsubsection{impl-dummy}

\code{DummyConsoleHandle}满足\code{Console}类型约束以便链接通过，但\code{write\_str}直接\code{unimplemented!()}，故意不做静默吞字，以便尽早暴露误用dummy feature的输出路径。仅适合无真实串口的编译/单测场景，生产内核应启用\code{impl-platform-console}。
